\section*{अध्याय \ref{ch:g11:cell-cycle-mitosis} --- कोशिका चक्र और समसूत्री विभाजन}

\begin{solution}{exo:g11:cell-cycle-mitosis:1}
$\mathrm{G_1}$: वृद्धि; S: \omterm{def:g11:dna-replication:replication}{DNA का प्रतिकृतिकरण}; $\mathrm{G_2}$: विभाजन
की तैयारी; M: \omterm{def:g11:cell-cycle-mitosis:cycle}{समसूत्री विभाजन} (\omterm{def:g10:cells-common-unit:organelle}{केंद्रक} का बँटना) और \omterm{def:g10:cells-common-unit:cell}{कोशिकाद्रव्य} विभाजन
(\omterm{def:g10:cells-common-unit:cell}{कोशिकाद्रव्य} का बँटना)।
\end{solution}

\begin{solution}{exo:g11:cell-cycle-mitosis:2}
\omterm{def:g11:cell-cycle-mitosis:chromosome}{क्रोमैटिड} प्रतिकृत \omterm{def:g11:cell-cycle-mitosis:chromosome}{गुणसूत्र} की दो एक जैसी प्रतियों में से एक है;
\omterm{def:g11:cell-cycle-mitosis:chromosome}{सेंट्रोमियर} वह क्षेत्र है जहाँ दोनों आपस में थमे रहते हैं। किसी \omterm{def:g11:cell-cycle-mitosis:chromosome}{गुणसूत्र}
में दो \omterm{def:g11:cell-cycle-mitosis:chromosome}{क्रोमैटिड} S के अंत से लेकर पश्चावस्था तक होते हैं।
\end{solution}

\begin{solution}{exo:g11:cell-cycle-mitosis:3}
पूर्वावस्था: \omterm{def:g11:cell-cycle-mitosis:chromosome}{गुणसूत्र} सिकुड़ते हैं, \omterm{def:g10:cells-common-unit:organelle}{केंद्रक} का लिफ़ाफ़ा टूटता है।
मध्यावस्था: \omterm{def:g11:cell-cycle-mitosis:chromosome}{गुणसूत्र} भूमध्य पर पंक्तिबद्ध होते हैं। पश्चावस्था:
\omterm{def:g11:cell-cycle-mitosis:chromosome}{क्रोमैटिड} अलग होकर ध्रुवों की ओर जाते हैं। अंत्यावस्था: \omterm{def:g10:cells-common-unit:organelle}{केंद्रक} के
लिफ़ाफ़े दोबारा बनते हैं, फिर \omterm{def:g10:cells-common-unit:cell}{कोशिकाद्रव्य} विभाजन।
\end{solution}

\begin{solution}{exo:g11:cell-cycle-mitosis:4}
S प्रावस्था में, प्रतिकृतिकरण के बीचोबीच।
\end{solution}

\begin{solution}{exo:g11:cell-cycle-mitosis:5}
46 \omterm{def:g11:cell-cycle-mitosis:chromosome}{गुणसूत्र}, 92 \omterm{def:g11:cell-cycle-mitosis:chromosome}{क्रोमैटिड}।
\end{solution}

\begin{solution}{exo:g11:cell-cycle-mitosis:6}
$\mathrm{G_1}$ \qty{11}{h}, S \qty{8}{h}, $\mathrm{G_2}$ \qty{4}{h},
M \qty{1}{h}। $2q$ वाली \omterm{def:g10:cells-common-unit:cell}{कोशिका} $\mathrm{G_2}$ में होती है, या
पश्चावस्था के अंत से पहले \omterm{def:g11:cell-cycle-mitosis:cycle}{समसूत्री विभाजन} में।
\end{solution}

\begin{solution}{exo:g11:cell-cycle-mitosis:7}
मध्यावस्था: 40 \omterm{def:g11:cell-cycle-mitosis:chromosome}{गुणसूत्र}, 80 \omterm{def:g11:cell-cycle-mitosis:chromosome}{क्रोमैटिड}। पश्चावस्था: 80 \omterm{def:g11:cell-cycle-mitosis:chromosome}{गुणसूत्र} (एक-एक
\omterm{def:g11:cell-cycle-mitosis:chromosome}{क्रोमैटिड} वाले), हर ध्रुव पर 40. पुत्री: 40 \omterm{def:g11:cell-cycle-mitosis:chromosome}{गुणसूत्र}, 40 \omterm{def:g11:cell-cycle-mitosis:chromosome}{क्रोमैटिड}।
\end{solution}

\begin{solution}{exo:g11:cell-cycle-mitosis:8}
समसूत्री सूचकांक $80/800 = 10\%$; \omterm{def:g11:cell-cycle-mitosis:cycle}{समसूत्री विभाजन} \qty{2.4}{h}।
पूर्वावस्था $48/800 \times 24 = \qty{1.44}{h}$; मध्यावस्था \qty{0.48}{h}
(29 मिनट); पश्चावस्था \qty{0.18}{h} (11 मिनट); अंत्यावस्था \qty{0.3}{h}
(18 मिनट)।
\end{solution}

\begin{solution}{exo:g11:cell-cycle-mitosis:9}
तर्कु के बिना \omterm{def:g11:cell-cycle-mitosis:chromosome}{क्रोमैटिड} अलग खींचे ही नहीं जा सकते: \omterm{def:g10:cells-common-unit:cell}{कोशिका} मध्यावस्था में
अटक जाती है, अपने सिकुड़े और पूरी तरह दिखते \omterm{def:g10:universal-dna:information}{गुणसूत्रों} के साथ, और हर एक
में दो \omterm{def:g11:cell-cycle-mitosis:chromosome}{क्रोमैटिड}। \omterm{def:g11:cell-cycle-mitosis:chromosome}{गुणसूत्र-चित्र} के लिए \omterm{def:g10:universal-dna:information}{गुणसूत्रों} की तस्वीर ठीक इसी
अवस्था में ली जाती है।
\end{solution}

\begin{solution}{exo:g11:cell-cycle-mitosis:10}
दो मीटर के उलझे धागे को छाँटा नहीं जा सकता; पर कुछ माइक्रोमीटर लंबी
छड़ों में सिकुड़ जाने पर वे 46 \omterm{def:g11:cell-cycle-mitosis:chromosome}{गुणसूत्र} बिना टूटे और बिना उलझे पंक्ति में
लगाए और अलग खींचे जा सकते हैं। बाद में वे इसलिए खुल जाते हैं कि \omterm{def:g10:universal-dna:gene}{जीन} केवल
ढीले रूप से ही पढ़े जा सकते हैं।
\end{solution}

\begin{solution}{exo:g11:cell-cycle-mitosis:11}
दस विभाजन: $2^{10} = 1024$ \omterm{def:g10:cells-common-unit:cell}{कोशिकाएँ}। नहीं: शुरुआती भ्रूण विभाजनों के बीच
बढ़ता नहीं, और उसके चक्र लगभग पूरे S तथा M ही होते हैं, $\mathrm{G_1}$
मुश्किल से ही।
\end{solution}

\begin{solution}{exo:g11:cell-cycle-mitosis:12}
एक पुत्री में 47 \omterm{def:g11:cell-cycle-mitosis:chromosome}{गुणसूत्र} होते हैं (एक अतिरिक्त), दूसरी में 45 (एक कम);
और उनका \omterm{def:g10:universal-dna:information}{DNA} $q$ में एक \omterm{def:g11:cell-cycle-mitosis:chromosome}{गुणसूत्र} भर जोड़ या घटाकर बनता है। किसी शरीर
\omterm{def:g10:cells-common-unit:cell}{कोशिका} में यह ग़लती लाखों करोड़ों में से एक \omterm{def:g10:cells-common-unit:cell}{कोशिका} पर असर डालती है, जो
आमतौर पर मर जाती है या हटा दी जाती है; पर यदि वह \omterm{def:g10:cells-common-unit:cell}{कोशिका} बची रहे और
विभाजित हो, और खोया या पाया गया \omterm{def:g11:cell-cycle-mitosis:chromosome}{गुणसूत्र} उसकी वृद्धि का नियंत्रण बिगाड़
दे, तो कोई अर्बुद शुरू हो सकता है।
\end{solution}

\begin{solution}{exo:g11:cell-cycle-mitosis:13}
प्रति दिन लगभग $\num{2e11}$ विभाजन, यानी किसी भी क्षण $\num{2e11}$
\omterm{def:g10:cells-common-unit:cell}{कोशिकाएँ} चक्र में (हर एक दिन भर में एक नई \omterm{def:g10:cells-common-unit:cell}{कोशिका} देती है)। \omterm{def:g11:cell-cycle-mitosis:cycle}{समसूत्री
विभाजन} रोक देने से आपूर्ति रुक जाती है; लाल \omterm{def:g10:cells-common-unit:cell}{कोशिकाएँ} 120 दिन जीती हैं पर
कुछ श्वेत \omterm{def:g10:cells-common-unit:cell}{कोशिकाएँ} केवल दिन भर, इसलिए हफ़्ते भर में श्वेत \omterm{def:g10:cells-common-unit:cell}{कोशिकाओं} की
गिनती ढह जाती है और संक्रमण पीछे-पीछे आते हैं --- यही कैंसररोधी दवाओं का
जाना-पहचाना दुष्प्रभाव है।
\end{solution}

\begin{solution}{exo:g11:cell-cycle-mitosis:14}
अंतरावस्था में \omterm{def:g11:cell-cycle-mitosis:chromosome}{गुणसूत्र} खुले और अदृश्य होते हैं; और पश्चावस्था में वे चल
रहे तथा आपस में मिले हुए होते हैं। मध्यावस्था में वे पूरी तरह सिकुड़े,
अलग-अलग, अब भी दो \omterm{def:g11:cell-cycle-mitosis:chromosome}{क्रोमैटिडों} के बने और एक ही तल में पड़े होते हैं: हर एक
को उसके आकार और रूप से पहचाना जा सकता है।
\end{solution}

\begin{solution}{exo:g11:cell-cycle-mitosis:15}
$\mathrm{G_1}$ में विश्राम कर रही यकृत \omterm{def:g10:cells-common-unit:cell}{कोशिकाएँ} चक्र में लौट आईं, और
द्रव्यमान बहाल होने तक बार-बार S, $\mathrm{G_2}$ तथा \omterm{def:g11:cell-cycle-mitosis:cycle}{समसूत्री विभाजन} से
गुज़रीं, फिर विश्राम में लौट गईं। दिन 0 पर लगभग सारी \omterm{def:g10:cells-common-unit:cell}{कोशिकाएँ} $q$ थामे
थीं; और दिन 3 पर बहुत-सी $q$ और $2q$ के बीच (S में) या $2q$ पर
($\mathrm{G_2}$ या \omterm{def:g11:cell-cycle-mitosis:cycle}{समसूत्री विभाजन} में)।
\end{solution}

\begin{solution}{pb:g11:cell-cycle-mitosis:1}
\textbf{1.} $100/1000 = 10\%$।

\textbf{2.} $0.10 \times 20 = \qty{2}{h}$।

\textbf{3.} पूर्वावस्था $60/1000 \times 20 = \qty{1.2}{h}$ (72 मिनट);
मध्यावस्था $0.36$ घंटे (22 मिनट); पश्चावस्था $0.16$ घंटे (लगभग 10 मिनट);
अंत्यावस्था $0.28$ घंटे (17 मिनट)।

\textbf{4.} पश्चावस्था। \omterm{def:g11:cell-cycle-mitosis:chromosome}{क्रोमैटिड} पहले ही सिकुड़े और जुड़े हुए हैं; और वह
हलचल कुछ माइक्रोमीटर की एक खींच भर है, जिसे मिनटों में छोटे होते तंतु कर
देते हैं।

\textbf{5.} विरली अवस्थाएँ कुछ ही \omterm{def:g10:cells-common-unit:cell}{कोशिकाओं} में दिखती हैं, और छोटी गिनती
बड़ी त्रुटि देती है; और चक्र केवल वृद्धि क्षेत्र में चलता है --- कहीं और
सूचकांक शून्य होता और आँकड़ों को पतला कर देता।

\textbf{6.} 16 \omterm{def:g11:cell-cycle-mitosis:chromosome}{गुणसूत्र}, 32 \omterm{def:g11:cell-cycle-mitosis:chromosome}{क्रोमैटिड}।

\textbf{7.} 32 \omterm{def:g11:cell-cycle-mitosis:chromosome}{गुणसूत्र} चल रहे होते हैं, हर ध्रुव तक 16.

\textbf{8.} पूर्वावस्था $2q$; पश्चावस्था में पूरी \omterm{def:g10:cells-common-unit:cell}{कोशिका} में $2q$; और
अंत्यावस्था में प्रति \omterm{def:g10:cells-common-unit:organelle}{केंद्रक} $q$।

\textbf{9.} S अंतरावस्था का $7/19$ है: यानी लगभग $900 \times 7/19 \approx
330$ \omterm{def:g10:cells-common-unit:cell}{कोशिकाएँ}।

\textbf{10.} पश्चावस्था: \omterm{def:g11:cell-cycle-mitosis:chromosome}{सेंट्रोमियर} अभी-अभी बँटे हैं और 16 \omterm{def:g10:universal-dna:information}{गुणसूत्रों} के
32 \omterm{def:g11:cell-cycle-mitosis:chromosome}{क्रोमैटिड} अब 32 \omterm{def:g11:cell-cycle-mitosis:chromosome}{गुणसूत्र} बनकर ध्रुवों की ओर जा रहे हैं।

\textbf{11.} प्रति घंटा $\num{20000}/20 = 1000$ नई \omterm{def:g10:cells-common-unit:cell}{कोशिकाएँ}।

\textbf{12.} पूरे क्षेत्र में प्रति घंटा 1000 \omterm{def:g10:cells-common-unit:cell}{कोशिकाएँ}; और एक क़तार में
वह क्षेत्र हर दिशा में शायद सौ \omterm{def:g10:cells-common-unit:cell}{कोशिकाएँ} चौड़ा है, इसलिए मोटे तौर पर प्रति
क़तार प्रति दिन $\num{24000}/(100 \times 100) \approx 2$--3 \omterm{def:g10:cells-common-unit:cell}{कोशिकाएँ}, और
हर एक \qty{100}{\micro m} की: यानी मिलीमीटर के कुछ दसवें हिस्से। देखी गई
\qty{1}{mm} का मतलब है कि वह क्षेत्र मानी हुई चौड़ाई से सँकरा है या
\omterm{def:g10:cells-common-unit:cell}{कोशिकाएँ} उससे अधिक लंबी होती हैं; परिमाण की कोटि सही है।

\textbf{13.} एक प्रारंभ स्थल $2 \times 50 \times 7 \times 3600 =
\num{2.5e6}$ क्षारक-जोड़े नक़ल करता है; इसलिए
$\num{1.6e10}/\num{2.5e6} \approx 6400$ प्रारंभ स्थल कम से कम।

\textbf{14.} हर अवस्था अधिक देर चलती है, इसलिए प्रति दिन कम \omterm{def:g10:cells-common-unit:cell}{कोशिकाएँ} चक्र
पूरा करती हैं; और यदि सारी प्रावस्थाएँ बराबर धीमी हों तो अनुपात, और इसलिए
समसूत्री सूचकांक, नहीं बदलते।

\textbf{15.} जो \omterm{def:g10:cells-common-unit:cell}{कोशिकाएँ} S पार कर चुकी थीं वे पहले घंटों में
$\mathrm{G_2}$ और \omterm{def:g11:cell-cycle-mitosis:cycle}{समसूत्री विभाजन} पूरा कर लेतीं; और दिन भर बाद कोई भी
\omterm{def:g10:cells-common-unit:cell}{कोशिका} \omterm{def:g11:cell-cycle-mitosis:cycle}{समसूत्री विभाजन} तक नहीं पहुँच सकती, इसलिए पूर्वावस्था से
अंत्यावस्था तक सब ग़ायब हो चुकी होतीं और केवल अंतरावस्था वाली \omterm{def:g10:cells-common-unit:cell}{कोशिकाएँ}
बचतीं, जो $\mathrm{G_1}$ में या S के आरंभ में अटकी होतीं।

\textbf{16.} $4/1000 \times 20 = \qty{0.08}{h}$, यानी लगभग 5 मिनट, 10 के
मुक़ाबले। चौंकाने वाली बात नहीं: 4 या 8 \omterm{def:g10:cells-common-unit:cell}{कोशिकाओं} के साथ गिनती पर संयोग ही
हावी रहता है।

\textbf{17.} $12/2000 \times 20 = \qty{0.12}{h}$, यानी लगभग 7 मिनट।

\textbf{18.} मानव: $0.04 \times 24 \approx \qty{1}{h}$; लहसुन
\qty{2}{h}। कोटि मिलती-जुलती है, और पादप \omterm{def:g10:cells-common-unit:cell}{कोशिकाएँ} कुछ धीमी हैं।

\textbf{19.} अर्बुद की \omterm{def:g10:cells-common-unit:cell}{कोशिकाएँ} लगातार चक्र चलाती हैं जबकि अधिकांश
सामान्य \omterm{def:g10:cells-common-unit:cell}{कोशिकाएँ} $\mathrm{G_1}$ में विश्राम करती हैं, इसलिए किसी भी क्षण
एक बड़ा अंश \omterm{def:g11:cell-cycle-mitosis:cycle}{समसूत्री विभाजन} में होता है। इसलिए तर्कु रोकने वाली दवा
अर्बुद की \omterm{def:g10:cells-common-unit:cell}{कोशिकाओं} पर सामान्य \omterm{def:g10:cells-common-unit:cell}{कोशिकाओं} से अधिक बार पड़ती है --- पर उन
सामान्य ऊतकों पर भी जो तेज़ी से विभाजित होते हैं।

\textbf{20.} चक्र 20 घंटे: $\mathrm{G_1}$ 8, S 7, $\mathrm{G_2}$ 4, और
\omterm{def:g11:cell-cycle-mitosis:cycle}{समसूत्री विभाजन} लगभग 2 (पूर्वावस्था 72 मिनट, मध्यावस्था 22, पश्चावस्था
7--10, अंत्यावस्था 17)। पश्चावस्था, जो सबसे छोटी है, वही अवस्था है जब
\omterm{def:g11:cell-cycle-mitosis:chromosome}{सेंट्रोमियर} बँटते हैं और हर ध्रुव को सोलहों \omterm{def:g10:universal-dna:information}{गुणसूत्रों} में से हर एक का एक
\omterm{def:g11:cell-cycle-mitosis:chromosome}{क्रोमैटिड} मिलता है।
\end{solution}
